% GENERATED FILE. Edit canonical lessons, then run npm run generate:books.
% canonical-source-hash: fnv1a64:049799f861119078
% canonical-lessons: MR-C396-aggadi, MR-C396-rul, MR-C396-lohmarga, MR-C396-malgadi, MR-C396-khadkhadne

\chapter{Trains and Rails}
\label{ch:mr-trains-and-rails}
\hlchaptermodality{396}

\begin{chapteropening}
\textbf{By the end of this chapter:} \emph{I can talk about trains, rails, railway lines and goods trains.}

Complete the last lesson of chapter 396: I can talk about trains, rails, railway lines and goods trains.
\end{chapteropening}

\section[\texorpdfstring{\mr{आगगाडी} (āggāḍī)}{आगगाडी (āggāḍī)}]{\mr{आगगाडी} (āggāḍī) — a train; a steam train}

\label{lesson:MR-C396-aggadi}



\begin{quote}
Before the new one: say the Marathi for dance, then the Marathi for a stage.
\end{quote}

\subsection*{You'll want to know: \mr{आगगाडी}}
\textbf{\mr{आगगाडी}} — \emph{āggāḍī} — \textquotedblleft{}a train; a steam train\textquotedblright{}.

\begin{quote}
\mr{आगगाडी} \mr{पुलावरून} \mr{हळू} \mr{जाते}.
\end{quote}

\begin{quote}
\emph{āggāḍī pulāvarūn haḷū jāte.}
\end{quote}

\begin{quote}
The train goes slowly over the bridge.
\end{quote}

Feminine. Plural: \mr{आगगाड्या}, \emph{āggāḍyā}.

The plural is \mr{आगगाड्या}; in everyday talk \mr{गाडी} alone often means the train.

\begin{cousinweb}
\textbf{\mr{आगगाडी}} (\emph{āggāḍī}) is a Marathi compound from the age of steam: \textbf{\mr{आग}}, \textquotedblleft{}fire\textquotedblright{}, inherited from Sanskrit \emph{agni} (CDIAL 55), and \textbf{\mr{गाडी}}, \textquotedblleft{}a cart\textquotedblright{}, inherited from \emph{*gāḍḍa} (CDIAL 4116).
\end{cousinweb}

\begin{grammarlens}[title={What you've built}]
One more word for getting around.
\end{grammarlens}

\subsection*{Your turn}
\begin{itemize}
\raggedright
  \item \emph{Say it:} \emph{āggāḍī}
  \item \emph{Say it:} \emph{āggāḍī}, once more
  \item \emph{Say it:} \emph{āggāḍī pulāvarūn haḷū jāte.} — The train goes slowly over the bridge.
  \item \emph{Recall:} say the Marathi for dance, then the Marathi for a stage, then say \emph{āggāḍī} again
\end{itemize}

\begin{tcolorbox}[breakable,colback=teal!4,colframe=teal!35!black,title={Before you move on}]
What does \mr{आगगाडी} mean? (\textquotedblleft{}A train; a steam train\textquotedblright{}.) Say it once more.
\end{tcolorbox}
\section[\texorpdfstring{\mr{रूळ} (rūḷ)}{रूळ (rūḷ)}]{\mr{रूळ} (rūḷ) — a rail; a railway track}

\label{lesson:MR-C396-rul}



\begin{quote}
Before the new one: say the Marathi for a stage, then the Marathi for a train.
\end{quote}

\subsection*{You'll want to know: \mr{रूळ}}
\textbf{\mr{रूळ}} — \emph{rūḷ} — \textquotedblleft{}a rail; a railway track\textquotedblright{}.

\begin{quote}
\mr{आगगाडी} \mr{रुळावरून} \mr{घसरली}.
\end{quote}

\begin{quote}
\emph{āggāḍī ruḷāvarūn ghasarlī.}
\end{quote}

\begin{quote}
The train came off the rails.
\end{quote}

Masculine.

Before an ending the long \mr{ऊ} shortens: \mr{रूळ}, but \mr{रुळावर}, on the rails.

\begin{cousinweb}
\textbf{\mr{रूळ}} (\emph{rūḷ}) names one steel rail of a track, and in the plural the line itself; where the word comes from is not settled.
\end{cousinweb}

\begin{grammarlens}[title={What you've built}]
One more word for getting around.
\end{grammarlens}

\subsection*{Your turn}
\begin{itemize}
\raggedright
  \item \emph{Say it:} \emph{rūḷ}
  \item \emph{Say it:} \emph{rūḷ}, once more
  \item \emph{Say it:} \emph{āggāḍī ruḷāvarūn ghasarlī.} — The train came off the rails.
  \item \emph{Recall:} say the Marathi for a stage, then the Marathi for a train, then say \emph{rūḷ} again
\end{itemize}

\begin{tcolorbox}[breakable,colback=teal!4,colframe=teal!35!black,title={Before you move on}]
What does \mr{रूळ} mean? (\textquotedblleft{}A rail; a railway track\textquotedblright{}.) Say it once more.
\end{tcolorbox}
\section[\texorpdfstring{\mr{लोहमार्ग} (lohmārga)}{लोहमार्ग (lohmārga)}]{\mr{लोहमार्ग} (lohmārga) — a railway line}

\label{lesson:MR-C396-lohmarga}



\begin{quote}
Before the new one: say the Marathi for a train, then the Marathi for a rail.
\end{quote}

\subsection*{You'll want to know: \mr{लोहमार्ग}}
\textbf{\mr{लोहमार्ग}} — \emph{lohmārga} — \textquotedblleft{}a railway line\textquotedblright{}.

\begin{quote}
\mr{हा} \mr{लोहमार्ग} \mr{डोंगरातून} \mr{जातो}.
\end{quote}

\begin{quote}
\emph{hā lohmārga ḍoṅgarātūn jāto.}
\end{quote}

\begin{quote}
This railway line runs through the hills.
\end{quote}

Masculine.

A formal word, common on signs and in the news; people say \mr{रेल्वे} or \mr{गाडीचा} \mr{रस्ता}.

\begin{cousinweb}
\textbf{\mr{लोहमार्ग}} (\emph{lohmārga}) is built from two tatsama words, Sanskrit \emph{loha}, \textquotedblleft{}iron\textquotedblright{}, and \emph{mārga}, \textquotedblleft{}road, way\textquotedblright{} (Monier-Williams): the iron road.
\end{cousinweb}

\begin{grammarlens}[title={What you've built}]
One more word for getting around.
\end{grammarlens}

\subsection*{Your turn}
\begin{itemize}
\raggedright
  \item \emph{Say it:} \emph{lohmārga}
  \item \emph{Say it:} \emph{lohmārga}, once more
  \item \emph{Say it:} \emph{hā lohmārga ḍoṅgarātūn jāto.} — This railway line runs through the hills.
  \item \emph{Recall:} say the Marathi for a train, then the Marathi for a rail, then say \emph{lohmārga} again
\end{itemize}

\begin{tcolorbox}[breakable,colback=teal!4,colframe=teal!35!black,title={Before you move on}]
What does \mr{लोहमार्ग} mean? (\textquotedblleft{}A railway line\textquotedblright{}.) Say it once more.
\end{tcolorbox}
\section[\texorpdfstring{\mr{मालगाडी} (mālgāḍī)}{मालगाडी (mālgāḍī)}]{\mr{मालगाडी} (mālgāḍī) — a goods train}

\label{lesson:MR-C396-malgadi}



\begin{quote}
Before the new one: say the Marathi for a rail, then the Marathi for a railway line.
\end{quote}

\subsection*{You'll want to know: \mr{मालगाडी}}
\textbf{\mr{मालगाडी}} — \emph{mālgāḍī} — \textquotedblleft{}a goods train\textquotedblright{}.

\begin{quote}
\mr{रात्री} \mr{एक} \mr{लांब} \mr{मालगाडी} \mr{गेली}.
\end{quote}

\begin{quote}
\emph{rātrī ek lāmb mālgāḍī gelī.}
\end{quote}

\begin{quote}
A long goods train went past at night.
\end{quote}

Feminine.

It carries freight, not people; a passenger train is a \mr{प्रवासी} \mr{गाडी}.

\begin{cousinweb}
\textbf{\mr{मालगाडी}} (\emph{mālgāḍī}) is a Marathi compound of \textbf{\mr{माल}}, \textquotedblleft{}goods\textquotedblright{}, a loan usually traced to Arabic \emph{māl}, \textquotedblleft{}property, goods\textquotedblright{}, and \textbf{\mr{गाडी}}, \textquotedblleft{}a cart, a train\textquotedblright{} (CDIAL 4116).
\end{cousinweb}

\begin{grammarlens}[title={What you've built}]
One more word for getting around.
\end{grammarlens}

\subsection*{Your turn}
\begin{itemize}
\raggedright
  \item \emph{Say it:} \emph{mālgāḍī}
  \item \emph{Say it:} \emph{mālgāḍī}, once more
  \item \emph{Say it:} \emph{rātrī ek lāmb mālgāḍī gelī.} — A long goods train went past at night.
  \item \emph{Recall:} say the Marathi for a rail, then the Marathi for a railway line, then say \emph{mālgāḍī} again
\end{itemize}

\begin{tcolorbox}[breakable,colback=teal!4,colframe=teal!35!black,title={Before you move on}]
What does \mr{मालगाडी} mean? (\textquotedblleft{}A goods train\textquotedblright{}.) Say it once more.
\end{tcolorbox}
\section[\texorpdfstring{\mr{खडखडणे} (khaḍkhaḍṇe)}{खडखडणे (khaḍkhaḍṇe)}]{\mr{खडखडणे} (khaḍkhaḍṇe) — to rattle}

\label{lesson:MR-C396-khadkhadne}



\begin{quote}
Before the new one: say the Marathi for a railway line, then the Marathi for a goods train.
\end{quote}

\subsection*{You'll want to know: \mr{खडखडणे}}
\textbf{\mr{खडखडणे}} — \emph{khaḍkhaḍṇe} — \textquotedblleft{}to rattle\textquotedblright{}.

\begin{quote}
\mr{जुना} \mr{डबा} \mr{रुळावर} \mr{खडखडत} \mr{होता}.
\end{quote}

\begin{quote}
\emph{junā ḍabā ruḷāvar khaḍkhaḍat hotā.}
\end{quote}

\begin{quote}
The old coach was rattling on the rails.
\end{quote}

Past (he …): \mr{तो} \mr{खडखडला}, \emph{to khaḍkhaḍlā}.

An echo word made into a verb: \mr{खड} \mr{खड} is the noise itself.

\begin{cousinweb}
\textbf{\mr{खडखडणे}} (\emph{khaḍkhaḍṇe}) is inherited and imitates its sound: CDIAL lists Marathi \emph{khaḍkhaḍṇẽ}, \textquotedblleft{}to rattle\textquotedblright{}, under Sanskrit \emph{khaṭakhaṭāyate}, \textquotedblleft{}crackles\textquotedblright{} (CDIAL 3771).
\end{cousinweb}

\begin{grammarlens}[title={What you've built}]
One more word for getting around.
\end{grammarlens}

\subsection*{Your turn}
\begin{itemize}
\raggedright
  \item \emph{Say it:} \emph{khaḍkhaḍṇe}
  \item \emph{Say it:} \emph{khaḍkhaḍṇe}, once more
  \item \emph{Say it:} \emph{junā ḍabā ruḷāvar khaḍkhaḍat hotā.} — The old coach was rattling on the rails.
  \item \emph{Recall:} say the Marathi for a railway line, then the Marathi for a goods train, then say \emph{khaḍkhaḍṇe} again
\end{itemize}

\begin{tcolorbox}[breakable,colback=teal!4,colframe=teal!35!black,title={Before you move on}]
What does \mr{खडखडणे} mean? (\textquotedblleft{}To rattle\textquotedblright{}.) Say it once more.
\end{tcolorbox}
